% Lösungen Einheit 3 von 5 – Zusammengesetzte Flächen, Maßstab und Volumen (zusammenbau v0.1)
\einheitenkopf{Einheit 3 von 5 · Zusammengesetzte Flächen, Maßstab und Volumen}

\erg{13}{a) Integral – der obere Rand ist gekrümmt \quad b) Integral $\frac{8}{3}$ plus Dreieck $4$; $A = \frac{20}{3} \approx 6{,}67$ \quad c) $3 + \frac{8}{3}$; $A = \frac{17}{3} \approx 5{,}67$ \quad d) $\frac{128}{3}$ Flächeneinheiten, eine Flächeneinheit $= 400\,\text{m}^2$; $A \approx 17\,066{,}7\,\text{m}^2 \approx 1{,}71$ ha \quad e) $A = \frac{8}{3}\,\text{dm}^2$; $V = 80\,\text{dm}^3$; Masse $= 192$ kg \quad f) $\frac{8}{3} + 2\pi$; $A \approx 8{,}95$ \quad g) (1) $y = 2x + 2$ ist die Sekante durch die Nullstelle $-1$ und den y-Achsenabschnitt; (2) Segment zwischen Graph und Sekante $= \frac{1}{6}$; (3) Dreieck unter der Sekante mit den Achsen; (4) das Segment ist ein Sechstel des Dreiecks.}
\erg{14}{a) Fehler: den Maßstab nur einfach angewandt. Richtig: eine Flächeneinheit $= 16\,\text{m}^2$, $A = 88\,\text{m}^2$ \quad b) Eine Flächeneinheit ist ein Quadrat mit einer Längeneinheit als Seite; beide Seiten werden mit dem Längenmaßstab gestreckt, der Inhalt also mit dessen Quadrat. \quad c) $A = \frac{32}{15} \approx 2{,}13\,\text{dm}^2$; $V = \frac{32}{15} \cdot 80 \approx 170{,}7\,\text{dm}^3$, also etwa $170{,}7$ Liter \quad d) Fläche unter dem Graphen von $0$ bis $2$ und rechts daran das Rechteck von $2$ bis $4$ mit der Höhe $f(2) = 3$.}

14d)

\begin{ksys}[xmin=-1,xmax=5,ymin=-1,ymax=5,klein]\flaeche{0.5*\x^2+1}{0}{2}{A_1} \flaeche{3}{2}{4}{A_2} \funktion{0.5*\x^2+1}{f}\end{ksys}

